\documentclass{beamer}

\usepackage[spanish]{babel}
\usepackage[utf8]{inputenc}
\usepackage[T1]{fontenc}
\usepackage{listings}


\usetheme{Madrid} % Puedes cambiar el tema según prefieras

\title{Gestión de Usuarios, Grupos y Permisos en Ubuntu}
\author{Tu Nombre}
\date{\today}

\begin{document}

%-------------------- PORTADA --------------------
\begin{frame}
    \titlepage
\end{frame}

%-------------------- ÍNDICE ---------------------
\begin{frame}{Contenido}
    \tableofcontents
\end{frame}

%=================================================
\section{Usuarios en Ubuntu}
%=================================================

%-------------------- Slide 1 --------------------
\begin{frame}{Conceptos Básicos}
    \begin{itemize}
        \item Cada usuario en Ubuntu tiene un \textbf{UID} (identificador único) y pertenece a uno o varios grupos.
        \item Los permisos de archivos y procesos dependen del usuario y de su(s) grupo(s).
    \end{itemize}
\end{frame}

%-------------------- Slide 2 --------------------
\begin{frame}{Comandos Básicos para la Gestión de Usuarios}
    \begin{center}
    \begin{tabular}{l l}
    \hline
    \textbf{Comando} & \textbf{Descripción} \\
    \hline
    whoami & Muestra el usuario actual \\
    who & Lista los usuarios conectados \\
    id & Muestra UID, GID y grupos del usuario \\
    adduser <usuario> & Crea un nuevo usuario (con home) \\
    userdel -r <usuario> & Elimina un usuario y su home \\
    passwd <usuario> & Cambia la contraseña de un usuario \\
    su <usuario> & Cambia al usuario especificado \\
    sudo <comando> & Ejecuta un comando como superusuario \\
    usermod -l <nuevo> <antiguo> & Cambia el nombre de un usuario \\
    \hline
    \end{tabular}
    \end{center}
\end{frame}

%-------------------- Slide 3 --------------------
\begin{frame}{Crear y Cambiar de Usuario}
    \textbf{Crear un usuario:}
    \begin{lstlisting}
sudo adduser juan
    \end{lstlisting}
    \textbf{Cambiar a otro usuario:}
    \begin{lstlisting}
su juan
exit
    \end{lstlisting}
\end{frame}

%=================================================
\section{Grupos en Ubuntu}
%=================================================

%-------------------- Slide 4 --------------------
\begin{frame}{Conceptos Básicos de Grupos}
    \begin{itemize}
        \item Un \textbf{grupo} en Ubuntu es un conjunto de usuarios con permisos compartidos.
        \item Útil para administrar accesos a archivos y directorios de forma colectiva.
    \end{itemize}
\end{frame}

%-------------------- Slide 5 --------------------
\begin{frame}{Comandos para Gestionar Grupos}
    \begin{center}
    \begin{tabular}{l l}
    \hline
    \textbf{Comando} & \textbf{Descripción} \\
    \hline
    groups <usuario> & Lista los grupos de un usuario \\
    groupadd <grupo> & Crea un nuevo grupo \\
    groupdel <grupo> & Elimina un grupo \\
    usermod -aG <grupo> <usuario> & Añade un usuario a un grupo \\
    gpasswd -d <usuario> <grupo> & Elimina un usuario de un grupo \\
    \hline
    \end{tabular}
    \end{center}
\end{frame}

%-------------------- Slide 6 --------------------

\begin{frame}[fragile]{Crear y Modificar Grupos}
    \textbf{Crear un grupo:}
    \begin{lstlisting}
sudo groupadd desarrolladores
    \end{lstlisting}
    \textbf{Agregar un usuario a un grupo:}
    \begin{lstlisting}
    sudo usermod -aG desarrolladores juan
    \end{lstlisting}
    \textbf{Eliminar un usuario de un grupo:}
    \begin{lstlisting}
sudo gpasswd -d juan desarrolladores
    \end{lstlisting}
\end{frame}

%=================================================
\section{Permisos en Ubuntu}
%=================================================

%-------------------- Slide 7 --------------------
\begin{frame}{Tipos de Permisos}
    \begin{itemize}
        \item \textbf{r (read)} = 4
        \item \textbf{w (write)} = 2
        \item \textbf{x (execute)} = 1
    \end{itemize}
    \vspace{0.3cm}
    Se aplican a:
    \begin{enumerate}
        \item \textbf{Usuario (Owner)}
        \item \textbf{Grupo (Group)}
        \item \textbf{Otros (Others)}
    \end{enumerate}
\end{frame}

%-------------------- Slide 8 --------------------
\begin{frame}{Ver Permisos de un Archivo}
    \begin{lstlisting}
ls -l archivo.txt
    \end{lstlisting}
    \begin{itemize}
        \item Ejemplo de salida: \texttt{-rw-r--r-- 1 juan desarrolladores ...}
        \item \textbf{Dueño:} juan
        \item \textbf{Grupo:} desarrolladores
        \item \textbf{Permisos:} Usuario = rw, Grupo = r, Otros = r
    \end{itemize}
\end{frame}

%=================================================
\section{Cambiar Permisos y Propietarios}
%=================================================

%-------------------- Slide 9 --------------------
\begin{frame}[fragile]{Cambiar Permisos con \texttt{chmod}}
    \textbf{Modo simbólico:}
    \begin{lstlisting}
chmod u+x script.sh    % Dar permiso de ejecución al usuario
chmod g-w archivo.txt  % Quitar permiso de escritura al grupo
    \end{lstlisting}
    \vspace{0.3cm}
    \textbf{Modo octal:}
    \begin{itemize}
        \item 7 (rwx), 6 (rw-), 5 (r-x), 4 (r--)
    \end{itemize}
    \begin{lstlisting}
chmod 755 script.sh
    \end{lstlisting}
\end{frame}

%-------------------- Slide 10 --------------------
\begin{frame}{Cambiar Propietario y Grupo}
    \textbf{Cambiar propietario (\texttt{chown}):}
    \begin{lstlisting}
sudo chown juan archivo.txt
sudo chown juan:desarrolladores archivo.txt
    \end{lstlisting}
    \vspace{0.3cm}
    \textbf{Cambiar grupo (\texttt{chgrp}):}
    \begin{lstlisting}
sudo chgrp desarrolladores archivo.txt
    \end{lstlisting}
\end{frame}

%=================================================
\section{Permisos Especiales}
%=================================================

%-------------------- Slide 11 --------------------
\begin{frame}{Sticky Bit (t)}
    \begin{itemize}
        \item Evita que otros usuarios borren archivos en directorios públicos.
        \item Comando:
        \begin{lstlisting}
chmod +t /directorio
        \end{lstlisting}
        \item Ejemplo: \texttt{/tmp} tiene este permiso por defecto.
    \end{itemize}
\end{frame}

%-------------------- Slide 12 --------------------
\begin{frame}{SetUID (s) y SetGID (s)}
    \textbf{SetUID (\texttt{u+s}):}
    \begin{itemize}
        \item Ejecuta un programa con los permisos de su propietario.
        \item \texttt{chmod u+s programa}
    \end{itemize}
    \vspace{0.3cm}
    \textbf{SetGID (\texttt{g+s}):}
    \begin{itemize}
        \item Los archivos creados heredan el grupo del directorio.
        \item \texttt{chmod g+s /directorio}
    \end{itemize}
\end{frame}

%=================================================
\section{Ejemplos Reales}
%=================================================

%-------------------- Slide 13 --------------------
\begin{frame}[fragile]{Administración de un Servidor Web}
    \begin{enumerate}
        \item Crear grupo \texttt{webdev}:
        \begin{lstlisting}
sudo groupadd webdev
        \end{lstlisting}
        \item Agregar usuarios al grupo:
        \begin{lstlisting}
sudo usermod -aG webdev juan
        \end{lstlisting}
        \item Asignar grupo y permisos en /var/www/html/:
        \begin{lstlisting}
sudo chown -R root:webdev /var/www/html/
sudo chmod -R 775 /var/www/html/
        \end{lstlisting}
    \end{enumerate}
\end{frame}

%-------------------- Slide 14 --------------------
\begin{frame}{Directorio Compartido}
    \begin{enumerate}
        \item Crear grupo \texttt{compartidos}:
        \begin{lstlisting}
sudo groupadd compartidos
        \end{lstlisting}
        \item Agregar usuarios al grupo:
        \begin{lstlisting}
sudo usermod -aG compartidos ana
        \end{lstlisting}
        \item Crear directorio y asignar permisos:
        \begin{lstlisting}
sudo mkdir /home/compartido
sudo chown -R root:compartidos /home/compartido
sudo chmod -R 770 /home/compartido
        \end{lstlisting}
    \end{enumerate}
\end{frame}

%=================================================
\section{Ejercicios Prácticos}
%=================================================

%-------------------- Slide 15 --------------------
\begin{frame}{Ejercicio 1: Creación de un Usuario}
    \textbf{Objetivo:} Crear un usuario \texttt{prueba1} y verificarlo.
    \begin{enumerate}
        \item \texttt{sudo adduser prueba1}
        \item Comprobar: \texttt{id prueba1}
    \end{enumerate}
    \textbf{Resultado esperado:} \texttt{uid=1001(prueba1) gid=1001(prueba1) ...}
\end{frame}

%-------------------- Slide 16 --------------------
\begin{frame}{Ejercicio 2: Creación de un Grupo y Asignación de Usuario}
    \begin{enumerate}
        \item \texttt{sudo groupadd practica}
        \item \texttt{sudo usermod -aG practica prueba1}
    \end{enumerate}
    \textbf{Comprobación:} \texttt{groups prueba1} \\
    \textbf{Esperado:} \texttt{prueba1 : prueba1 practica}
\end{frame}

%-------------------- Slide 17 --------------------
\begin{frame}{Ejercicio 3: Cambio de Propietario de un Archivo}
    \begin{enumerate}
        \item Crear \texttt{ejemplo.txt} en tu \textbf{home}.
        \item Cambiar propietario a \texttt{prueba1}.
    \end{enumerate}
    \textbf{Comprobación:} \texttt{ls -l ejemplo.txt} \\
    \textbf{Esperado:} \texttt{-rw-r--r-- 1 prueba1 usuario ... ejemplo.txt}
\end{frame}

%-------------------- Slide 18 --------------------
\begin{frame}{Ejercicio 4: Cambiar Permisos de un Archivo}
    \begin{enumerate}
        \item Crear \texttt{script.sh}.
        \item Dar permisos de ejecución a todos (\texttt{chmod 755 script.sh}).
    \end{enumerate}
    \textbf{Comprobación:} \texttt{ls -l script.sh} \\
    \textbf{Esperado:} \texttt{-rwxr-xr-x ... script.sh}
\end{frame}

%-------------------- Slide 19 --------------------
\begin{frame}{Ejercicio 5: Directorio Compartido con Grupo}
    \begin{enumerate}
        \item \texttt{sudo groupadd equipo1}
        \item Crear \texttt{usuario2} y añadirlo a \texttt{equipo1}.
        \item Directorio \texttt{/home/compartido/} con grupo \texttt{equipo1}.
        \item \texttt{chmod -R 770 /home/compartido/}
    \end{enumerate}
    \textbf{Comprobación:} \texttt{ls -ld /home/compartido/} \\
    \textbf{Esperado:} \texttt{drwxrwx--- ... root equipo1 ...}
\end{frame}

%-------------------- Slide 20 --------------------
\begin{frame}{Ejercicio 6: Configurar el Sticky Bit}
    \begin{enumerate}
        \item Crear \texttt{/home/publico/} con permisos de escritura para todos.
        \item Activar Sticky Bit: \texttt{chmod +t /home/publico/}
    \end{enumerate}
    \textbf{Comprobación:} \texttt{ls -ld /home/publico/} \\
    \textbf{Esperado:} \texttt{drwxrwxrwt 2 root root ... /home/publico/}
\end{frame}

%-------------------- Slide 21 --------------------
\begin{frame}{Ejercicio 7: Usuario con Shell Restringido}
    \begin{enumerate}
        \item Crear \texttt{limitado} con shell \texttt{/bin/rbash}.
        \item Probar \texttt{cd /} y ver error.
    \end{enumerate}
    \textbf{Comprobación:} \texttt{rbash: cd: restricted}
\end{frame}

%-------------------- Slide 22 --------------------
\begin{frame}{Ejercicio 8: Asignar Permisos con SetGID}
    \begin{enumerate}
        \item \texttt{sudo groupadd proyectos}
        \item Crear \texttt{/home/proyectos/} y asignar grupo.
        \item \texttt{chmod g+s /home/proyectos/}
    \end{enumerate}
    \textbf{Comprobación:} \texttt{ls -ld /home/proyectos/} \\
    \textbf{Esperado:} \texttt{drwxrws--- ... root proyectos ...}
\end{frame}

%=================================================
\section{Conclusiones}
%=================================================

%-------------------- Slide 23 --------------------
\begin{frame}{Conclusión}
    \begin{block}{Resumen}
    \begin{itemize}
        \item Administrar usuarios, grupos y permisos es esencial para la seguridad y eficiencia en Ubuntu.
        \item Comprender los \textbf{permisos (r, w, x)}, la asignación de propietarios y grupos permite un control preciso.
        \item Los ejercicios propuestos te ayudan a practicar y verificar tu aprendizaje de forma directa.
    \end{itemize}
    \end{block}
    \vspace{0.5cm}
    \centering
    \alert{¡Si tus resultados no coinciden con los esperados, revisa y repite los pasos!}
\end{frame}

%-------------------- FIN --------------------
\end{document}
